% titel: Blatt 0, A6: Quadratische und Wurzelgleichungen

% Teil I: Werkzeuge, kein Lösungsweg.

Quadratische Gleichungen und Wurzelgleichungen für die reelle Zahl $x$
lösen.

\begin{werkzeugkasten}[Quadratische Gleichung]
$ax^2+bx+c=0$ lösen per Mitternachtsformel
\begin{equation}
  x_{1,2}=\frac{-b\pm\sqrt{b^2-4ac}}{2a},
\end{equation}
oder nach Normierung $x^2+px+q=0$ per $p$-$q$-Formel
$x_{1,2}=-\frac p2\pm\sqrt{\left(\frac p2\right)^2-q}$.
Zwei, eine oder keine reelle Lösung je nach Diskriminante.
\end{werkzeugkasten}

\begin{werkzeugkasten}[Wurzelgleichung]
\begin{enumerate}
  \item Wurzel isolieren.
  \item Beide Seiten quadrieren (beim Quadrieren können Scheinlösungen
        entstehen).
  \item Die quadratische Gleichung lösen.
  \item \textbf{Probe in der Originalgleichung}: Quadrieren ist nicht
        äquivalent, negative Argumente der Wurzel und falsche Vorzeichen
        ausschließen.
\end{enumerate}
Bei zwei Wurzeln beide nacheinander isolieren bzw. einmal geeignet quadrieren
und die Wurzeln geschickt kombinieren.
\end{werkzeugkasten}

\begin{achtung}[Darauf achten]
\begin{itemize}
  \item Jede Wurzelgleichung-Lösung zurückprüfen; es können Lösungen
        auftreten, die die ursprüngliche Gleichung nicht erfüllen.
  \item Das Wurzelargument muss reell (nichtnegativ) sein.
  \item Vorzeichen: $10^{x^2}=100$ gibt $x=\pm\sqrt2$; beide Werte testen.
  \item In Gleichungen wie $\sqrt{x+1}=x-5$ ist die rechte Seite bei der
        Lösung $x=3$ negativ ($3-5=-2$) -- solche Scheinlösungen verwerfen.
\end{itemize}
\end{achtung}
